% =====================================================================
\chapter[结构收敛]{结构收敛：过约束解集与工程吸引子}\label{ch:convergence}
% =====================================================================

本章处理问题 \ref{prob:low}。高层结构确实收敛，但收敛机制与底层不同，不同强度的
命题必须分开陈述；收敛强度依赖领域绑定度，其载体是关系而非组件。本章末尾给出可
证伪的验证设计。

\section{三个强度不同的收敛命题}\label{sec:three-claims}

\begin{hypothesis}[弱命题：局部模式复现]\label{hyp:weak}
不同领域的软件反复使用同一批局部结构解法：分层与阶段化、机制与策略分离、控制与数据
分离、事件与状态机、声明式描述与执行引擎分离、前端与后端分离、资源所有权与生命周期
管理。这些结构的重要性不在于名称相同，而在于它们具有相同的\emph{结构性动机}。
\end{hypothesis}

弱命题有大量工程经验支持，其强度也最低：它几乎不可能被证伪，因此解释力有限。

\begin{hypothesis}[中命题：抽象后结构相似性]\label{hyp:medium}
对于来自不同领域、由不同团队开发、具有不同实现细节的大型软件系统，在消除实现语言、
组件命名与领域语义差异之后，其高层架构在某个可定义的抽象空间中呈现出\emph{显著高于
随机基线}的结构相似性。
\end{hypothesis}

中命题是可操作化的核心命题。它可被操作化（第 \ref{sec:baseline} 节），
且其成立与否不依赖任何形而上假设。

\begin{hypothesis}[强命题：单点收敛]\label{hyp:strong}
几乎所有正常偏置的大型软件最终都会收敛到同一种高层架构。
\end{hypothesis}

强命题证据不足，而且很可能过强。它的主要问题在于把"存在吸引子"与"只有一个吸引子"
混为一谈。

本文以式 \ref{hyp:medium}（中命题）为核心命题，把式 \ref{hyp:weak}（弱命题）视为其
退化形式，把式 \ref{hyp:strong}（强命题）视为过强表述而予以拒绝。

\section{约束系统与解集几何}\label{sec:overconstrain}

\begin{definition}[约束系统与解集]\label{def:constraint}
设设计空间为 $d$ 维流形 $\mathcal{D}$，一个设计问题由 $c$ 个独立约束
$g_1, \dots, g_c : \mathcal{D} \to \mathbb{R}$ 定义，其\term{解集}为
$\mathcal{S} = \{x \in \mathcal{D} : g_j(x) = 0,\; j = 1..c\}$。
\end{definition}

\begin{proposition}[过约束与欠约束：在一个正则性假设之下]\label{prop:overconstrain}
设 $\mathcal{S} \subseteq \mathcal{D}$ 被描述为映射 $g : \mathcal{D} \to
\mathbb{R}^{c}$ 的零点集，$0$ 是 $g$ 的正则值，且 $g$ 的 $c$ 个分量相互独立。则
\begin{enumerate}[label=(\alph*)]
\item 当 $c < d$ 时，若 $\mathcal{S} \ne \emptyset$，则它是维数为 $d - c$ 的子流形；
\item 当 $c = d$ 时，$\mathcal{S}$ 是离散的——零维的——并且仅在一个附加假设（例如
      $\mathcal{S}$ 的紧致性）之下才是有限的；
\item 当 $c > d$ 时，$\mathcal{S}$ 在一般位置下为空。
\end{enumerate}
\end{proposition}

\begin{proof}
在 $g$ 关于 $\{0\}$ 横截的条件下，$g$ 的每个独立分量把 $\mathcal{S}$ 的余维数降低
一，故由横截性定理（Thom）得 $\mathrm{codim}\,\mathcal{S} = c$。这在 $d - c > 0$ 时
给出维数 $d - c$；在 $c = d$ 时给出离散集合；而在 $c > d$ 时不可能成立，因为 $d$ 维
流形的子集不可能有超过 $d$ 的余维数——故原像在一般位置下为空。$c = d$ 时的有限性
不能由维数计数推出：非紧流形上的零维闭子集可以是无限的，$\mathcal{S}$（或
$\mathcal{D}$）的紧致性必须另行假设。
\end{proof}

对这一命题的读法需要两处更正，而这两处都比命题本身更重要。

第一，$c \ge d$ 的情形\emph{并不}给出"一般为有限或空集"。零维给出的是离散性，而离散
性只有在紧致性之下才给出有限性。更弱而正确的一般陈述是：当 $c \ge d$ 时解集没有内点
——它是低维的层（stratum），而不是一个开区域——这既是该论证所需的全部，也是它所能
支撑的全部。

第二，也是更实质的一点：这一前提假设很强，而所讨论的领域并不是该假设所描述的那类
对象。横截性加独立性是关于流形之间光滑映射的条件，而软件设计空间两者都不是：它是
离散的，其约束常常是不等式、枚举、兼容性关系以及彼此之间的依赖，其可行集通常是有限
或可数的候选集合，而不是流形。因此该命题是底层情形的一个\emph{模型}——一个有用的
模型，因为它解释了为什么增加约束会减少自由度，以及为什么底层显得被钉死——而不是
一条关于软件架构的定理。它应当被读作：\emph{在}设计空间容许光滑描述、且约束相互独立
\emph{之处}，自由度随约束数目线性收缩，并在 $c = d$ 时消失。

本章关于层级高低的实际主张在这次削弱之后完好无损，因为承载这一论证的从来不是维数
计数。经验主张是：在低层，约束源自已经存在的对象，因而彼此独立；而在高层，若干表面
上的约束只是同一个偏好的不同表述，因而无法满足独立性（附录 \ref{app:proofs} 给出了
同一点，并加以推广）。维数计数是对这一不对称性的一幅图像，而不是它的证明；这一不
对称性可检验的形式是假设 $H_2$，它问的是一个被测的领域绑定度是否单调地预测相似性
——这是一个关于数据的问题，而不是关于横截性的问题。

命题 \ref{prop:overconstrain} 给出"底层唯一、高层多解"的一个干净解释，但需要用词
谨慎。它\emph{不是}在说底层真的被形式化为一个约束系统，而是说底层设计空间的有效
维数很低、有效约束很多，因此其解集在几何意义上退化。

\begin{example}[底层]\label{ex:abi}
考虑为一个固定指令集架构定义函数调用约定。设计自由度包括参数的传递位置（寄存器或
栈）、寄存器的保存责任、栈帧布局、返回地址存放、可变参数的传递、对齐要求、以及
与既有二进制、调试器、异常展开表、系统调用接口和动态链接器的兼容性。这些约束中绝大
多数不是"可协商的偏好"，而是被\emph{已有对象}固定下来的：既有二进制必须继续工作，
调试器必须能回溯栈，异常必须能展开，内核必须能接收系统调用。因此 $c \ge d$，解集
退化为极少数可行点，在实践上常常退化为一。
\end{example}

\begin{example}[高层]\label{ex:arch}
考虑为一个任务处理系统选择架构。设计自由度包括：单体中的函数调用、多进程消息传递、
微服务 RPC、事件驱动异步处理、Actor 模型、分布式工作流引擎，以及它们在下列维度上的
连续参数：部署粒度、状态归属、一致性模型、同步或异步、失败边界、数据分区方式。
此时问题层面的约束显著少于自由度，$\mathcal{S}$ 是一个高维流形。
\end{example}

那么为什么高层仍然收敛？因为存在\emph{第二组约束}：它们不来自问题，而来自制造与
维护该系统的人类组织：

\begin{principle}[组织约束对解集的压缩]\label{pr:org}
问题层面欠约束的高维解流形，会被组织层面的约束压缩成少数几个尖峰。这些约束包括：
团队沟通成本、认知负荷上限、部署与运维能力边界、技术人才市场、既有技术栈的沉没
成本、以及监管与合规要求。
\end{principle}

由此得到一个推论：\emph{高层收敛不是朝向工程最优，而是朝向"有界组织可导航
的结构"}。这一点在后面的强度定律中会被精确化。

\section{收敛机制}\label{sec:mechanisms}

\begin{enumerate}[label=(M\arabic*)]
\item \term{复杂性约束}。设系统有 $N$ 个组件，若依赖不受组织，潜在有向依赖数达
      $N(N-1)$，无向达 $N(N-1)/2$。系统的可维护性因此要求对依赖施加结构：模块化、
      分层、显式接口、依赖方向约束、子系统边界。这些技术不完全相同，但都在解决同一个
      问题：\emph{限制复杂性传播}。模块化与信息隐藏的经典论证见 \cite{parnas1972}。
\item \term{信息流约束}。大量软件可被抽象为"输入 $\to$ 处理 $\to$ 状态变化 $\to$
      输出"（并叠加反馈、并发事件、外部副作用与持续运行的状态机）。当不同职责具有
      不同的变化频率、错误模式与资源需求时，分离它们具有实际收益：解析与执行分离、
      业务规则与数据访问分离、纯计算与副作用分离、控制平面与数据平面分离。
\item \term{状态管理约束}。系统必须回答：谁拥有某个状态？谁可以修改它？它何时创建和
      销毁？多个组件如何观察它的变化？并发访问如何协调？出错时如何恢复一致性？这些
      问题具有很强的普遍性，因而反复产生相似的结构：事务与锁管理、进程与句柄管理、
      GUI 的状态所有者与事件循环、复制状态机、实体生命周期与系统更新。
\item \term{认知约束}。大型软件不仅受机器约束，也受开发者的理解能力约束。若所有组件
      之间都存在隐式关系，则即使系统在计算意义上可以运行，开发者也无法预测修改的影响。
      因此人倾向于构建层次、命名空间、模块、明确接口、可组合抽象与局部可理解的状态。
      \emph{软件架构不仅是在组织计算，也是在有限认知能力下组织可理解性。}
\item \term{演化约束}。大型软件不是一次性设计完成的，而经历初始实现、功能增长、性能
      优化、故障修复、团队交接、模块重构与兼容性维护。某些决策会随演化逐渐显露成本，
      于是产生筛选效应。现实中的演化不保证最优，历史包袱、兼容性要求、
      组织结构与沉没成本都可能让不理想的架构长期存活。
\end{enumerate}

机制 (M1)--(M3) 来自问题的客观结构，机制 (M4)--(M5) 来自人类与历史。这一区分对应
第 \ref{sec:strength} 节的强度定律：前三种机制产生跨领域相似性，后两种产生领域内
相似性与路径依赖。

机制 (M4) 与 (M5) 的经典表述分别是认知负荷下的层级分解\cite{simon1962}与沟通结构
对系统结构的决定作用\cite{conway1968}。"无终态"架构观与多稳定形态之间的振荡见
\cite{nygard2018}；组织形态通过团队拓扑反向驱动架构的实践见限界上下文与领域驱动
设计\cite{evans2003}。

\section{收敛强度与领域绑定度}\label{sec:strength}

\begin{definition}[领域绑定度]\label{def:dom}
设层级 $\Lambda$ 的语义依赖一组领域公理 $\mathrm{Ax}(\Lambda)$（例如"账户余额可以为
负"或"帧必须在 16.6\,ms 内提交"）。定义 $\Lambda$ 的\term{领域绑定度}为
\begin{equation}\label{eq:dom}
\mathrm{Dom}(\Lambda) \;=\; \frac{\bigl|\mathrm{Ax}(\Lambda)\bigr|}
{\bigl|\mathrm{Ax}(\Lambda)\bigr| + \bigl|\mathrm{Ax}_{\text{通用}}(\Lambda)\bigr|} ,
\end{equation}
即领域特有公理在该层全部公理中的占比。$\mathrm{Dom} = 0$ 表示该层完全由通用公理
决定（如内存局部性、并发协调成本），$\mathrm{Dom} = 1$ 表示该层完全是领域约定。
\end{definition}

\begin{conjecture}[收敛强度定律]\label{conj:strength}
设 $\mathrm{Conv}(\Lambda)$ 为层级 $\Lambda$ 在跨系统比较中的结构相似性度量
（第 \ref{sec:baseline} 节给出操作化）。则 $\mathrm{Conv}$ 关于
$\mathrm{Dom}$ 单调递减：
\begin{equation}\label{eq:convlaw}
\Lambda_1 \preceq \Lambda_2 \;\Longrightarrow\;
\mathrm{Conv}(\Lambda_1) \;\ge\; \mathrm{Conv}(\Lambda_2)
\qquad \text{当} \qquad
\mathrm{Dom}(\Lambda_1) \le \mathrm{Dom}(\Lambda_2).
\end{equation}
\end{conjecture}

猜想 \ref{conj:strength} 有可观察的定性支持：越靠近物理与网络层，趋同越强（分布式
计算的八条谬误是无协商余地的硬约束）；越靠近业务域，趋同越弱（保险核心系统与推荐
系统在领域模型上几乎不会同构）；越靠近组织结构，越受沟通结构支配；越靠近技术成熟期，
形态越收敛；越靠近核心与高风险模块，越保守趋同。

本文不把式 \eqref{eq:convlaw} 称为定律，因为它依赖 $\mathrm{Dom}$ 的可定义性与
$\mathrm{Conv}$ 的可测性，二者都需要在实验中被建立。它的价值是把一个模糊的方向性
判断转化为一个可被证伪的单调性假设。

\section{关系而非组件}\label{sec:relations-converge}

\begin{principle}[关系优先]\label{pr:relations}
跨系统的结构共性主要存在于\emph{关系}之中，而不是\emph{组件}之中。不同系统不一定
具有相同的模块，但可能具有相似的关系型：
\begin{itemize}
\item 哪些部分负责决策，哪些部分负责执行；
\item 哪些部分拥有状态，哪些部分只能读取状态；
\item 哪些部分定义稳定契约，哪些部分允许变化；
\item 哪些部分协调全局，哪些部分独立工作；
\item 哪些依赖可以跨越边界，哪些依赖必须被禁止。
\end{itemize}
\end{principle}

用定义 \ref{def:relgraph} 与 \ref{def:skeleton} 的语言：需要寻找的不是一个通用
的架构模板，而是一组能在不同系统中反复出现的结构关系及其\emph{成立条件}。这与
第 \ref{ch:infrastructure} 章的内核六职能是同一条思路的两个层次：六职能给的是
\emph{管理职责}的坐标，关系型给的是\emph{职责之间}的坐标。

\begin{example}[数据库与浏览器的抽象骨架]\label{ex:db-browser}
一个数据库引擎包含 SQL 解析器、查询优化器、执行引擎、缓冲池、事务管理器与存储引擎；
一个浏览器包含 HTML/CSS/JavaScript 解析器、样式与布局引擎、脚本执行引擎、渲染引擎、
网络子系统与缓存持久化机制。二者的模块名称、内部算法与具体职责都不同。但在适当的
骨架映射 $\phi$ 下，二者都表现为：\emph{输入与接口层}（接收请求、解析输入、建立执行
上下文）、\emph{解释与规划层}（把外部意图转换为内部可执行表示）、\emph{执行与协调层}
（调度工作、协调子系统、管理执行过程）、\emph{状态与资源层}（管理数据、缓存、生命
周期与持久化）。
\end{example}

例 \ref{ex:db-browser} 的价值有限：这个骨架已经宽泛到几乎可以描述任何
系统。因此需要一条对抽象层级的附加要求：

\begin{principle}[抽象层级的双向约束]\label{pr:abstraction-level}
有研究价值的抽象层级必须同时满足两个方向：\emph{向上}，它必须能识别跨系统的共同
结构；\emph{向下}，它必须保留足够的区分能力，使不同系统在该层级上仍然可被区分。
越抽象，实现差异越被消除，但描述也越容易失去区分力；存在一个恰当的中层。
\end{principle}

第 \ref{ch:infrastructure} 章的六职能坐标正是为满足这一双向要求而设计的：它比
"输入-处理-输出"窄（因此有区分力），比任何具体模块列表宽（因此跨领域可比）。

\section{验证设计}\label{sec:baseline}

中命题（假设 \ref{hyp:medium}）要成为科学命题，必须能回答"与什么比较"。一个完整的
检验设计包含五步。

\paragraph{第一步：建立架构图。}
为每个系统建立图 $\mathcal{G} = (V, E)$，节点代表组件或职责，边代表调用、依赖、
控制流或数据流。按定义 \ref{def:relgraph}，\emph{不同类型的关系应分开记录}，
避免把完全不同的关系混为一谈。

\paragraph{第二步：定义抽象映射。}
将具体组件映射到更一般的类别，例如：输入与解析、规划与决策、执行与协调、状态与
持久化、通信与边界、资源与生命周期管理。这些类别不能预先定义得过于宽泛，否则任何
系统都会显得相似。

\paragraph{第三步：选择有区分力的结构指标。}
可用的指标包括：依赖图的模块化程度、层次结构深度、跨模块依赖比例、状态所有权的
集中程度、控制流与数据流的耦合程度、核心组件的网络中心性、以及不同职责之间的依赖
方向。

\paragraph{第四步：构建合理的对照组。}
不能只比较几个著名软件然后发现它们都存在模块化。必须将观察到的相似性与适当的随机
基线比较，并控制系统规模、领域、平台、架构风格和历史技术栈等因素。

\paragraph{第五步：区分真正的收敛与共同来源。}
如果两个系统采用相同的架构模式，可能是因为：它们独立解决了相似问题；它们使用了相同
的框架；它们受到相同的标准影响；它们借鉴了相同的设计文献；它们继承了相同的历史架构。
\emph{只有在控制这些因素之后，才更有理由讨论独立的结构收敛。}

第 \ref{ch:program} 章把这一设计写成假设 $H_1$ 与 $H_2$。

\section{反例与边界}\label{sec:limits}

必须明确列出使收敛命题失效或退化的情况：

\begin{enumerate}[label=(\alph*)]
\item \term{路径依赖}。一旦某种结构被生态、工具链或标准锁定，它可能长期存留，即使
      次优。锁定的强度与迁移成本成正比，与被锁定结构的质量无关。
\item \term{规模相变}。在 $10^{2}$ 与 $10^{9}$ 个实体两种规模下，最优结构是定性不同
      的：一致性模型、分片策略、故障模型全部改变。因此不存在一个"大系统架构"，
      只存在"某个规模区间内的架构"。这直接限制了任何"所有大系统都收敛"的表述。
\item \term{组织政治}。沟通结构对系统结构的决定作用（Conway 定律）意味着，即使技术
      问题完全相同，不同组织仍会产生不同的系统边界。这既是收敛的部分原因（当组织
      形态趋同时），也是发散的部分原因（当组织形态不同时）。
\item \term{强耦合系统}。当系统内部耦合无法被局部化时（例如紧耦合的实时控制系统或
      单体内核），模块边界可能不是设计选择而是被迫的，此时比较"架构"本身的意义
      下降。
\item \term{抽象层级选择的人为性}。第三步与第四步的指标选择本身包含研究者自由度。
      不同的职责分类系统会给出不同的相似性结论。这是该路线最脆弱的方法论环节，
      必须通过预注册指标与多分类系统稳健性检验来缓解。
\end{enumerate}

本章的可证伪部分是假设 \ref{hyp:medium} 与猜想 \ref{conj:strength}：若在控制规模、
领域、平台与框架之后，跨系统结构相似性不显著高于随机基线，或相似性关于
$\mathrm{Dom}$ 的单调性不成立，则本章的核心主张被削弱（假设 $H_1$ 与 $H_2$，
第 \ref{ch:program} 章）。弱命题 \ref{hyp:weak} 与原则 \ref{pr:relations} 不受此
影响，因为二者是解释性框架而非经验断言。
